%%%%%%%%%%%%%%%%%%%%%%%%%%%%%%%%%%%%%%%%%%%%%%%%%%%%%%%%%%%%
%% CHAPTER 13 -- A Map of How LLMs Go Wrong
%%%%%%%%%%%%%%%%%%%%%%%%%%%%%%%%%%%%%%%%%%%%%%%%%%%%%%%%%%%%

\chapter{A Map of How LLMs Go Wrong}
\label{ntg:ch13}

\scenelabel
\begin{scenestrip}
\sceneitem{The problem}
\sceneitem{How we got here}
\sceneitem{Where we stand}
\sceneitem{What goes wrong}
\end{scenestrip}

\paragraph{The problem.} The twelve chapters that came before have each ended with a
section called \emph{What goes wrong, and why}. By now the
reader has met a substantial vocabulary of LLM failure modes:
hallucination, miscalibration, prompt brittleness, the
alignment tax, sycophancy, reward hacking, jailbreaks, lost-in-
the-middle, retrieval drift, citation hallucination, prompt
injection, confused-deputy attacks, cost runaways. Each of
these terms picks out a real and important kind of failure.
Each was paid for, in the chapter that introduced it, with a
careful explanation of the mechanism that produces it.

What the vocabulary does not give the reader, on its own, is a
way of saying to the colleague across the desk who has just
been handed a wrong answer by a language model and wants to
know, in plain words, what went wrong. ``The cross-entropy
objective trains for plausibility rather than truth, and the
retrieval pipeline appears to have surfaced a topically
adjacent but factually incorrect passage'' is correct. It is
not what the colleague came to hear. The problem of this final
chapter is the gap between the engineer's vocabulary, which is
precise, and the user's experience of error, which is
unmediated. We want a model of LLM failure that an intelligent
non-specialist can hold in their head, that maps onto the
technical mechanisms covered in previous chapters, and that
gives a user something to do when an answer arrives that they
suspect is wrong.

This chapter builds that model. The model has five axes,
arranged from the most mechanical to the most human. The
ordering is not arbitrary. It is the empirically observed order
of how reliably the model performs at each level: very well at
the bottom, less well as one ascends. Walking the ladder, when
an answer arrives that smells wrong, is the discipline this
chapter recommends.

\paragraph{How we got here.} The study of language as a layered phenomenon is older than
the language model. Linguists since at least the early
twentieth century have distinguished between the structure of
an utterance (its syntax), the meaning it conveys (its
semantics), and the situational effect it has (its
pragmatics). Each level has its own rules, its own
characteristic failures, and its own diagnostic vocabulary.

Noam Chomsky's distinction between competence and performance,
in the 1960s, separated the abstract rules of grammar from
the noisy reality of speech. A speaker can have competence in
a rule (they know it, in the sense that their judgements
about new sentences conform to it) without performing it
reliably (their actual speech is full of slips, restarts, and
errors that the rule does not predict). The distinction
matters for our purposes because it captures, almost
exactly, the gap between what the model has been trained to
do and what the model in any given moment actually does.

H. Paul Grice, in his 1967 William James lectures at Harvard,
articulated what he called the \emph{cooperative principle}:
that speakers and listeners in conversation are tacitly
cooperating to make the exchange go well, and that the
implications a listener draws from a speaker's utterance
depend on this assumption of cooperation. Grice gave us the
vocabulary of \emph{conversational implicature} --- the
distinction between what a speaker literally says and what
they mean by saying it --- and the four maxims of conversation
(quantity, quality, relation, manner) that participants are
assumed to follow. None of this was about machines. All of it
turns out to matter when the machine begins to talk.

John Searle's \emph{speech acts} of 1969 categorised what
utterances do as well as what they say. A statement asserts; a
question requests information; a promise commits the speaker
to a future action. The taxonomy gave us the vocabulary in
which to ask, of any utterance, not only ``is it true'' but
``is it the right kind of thing to be saying here''. This
question turns out to be unexpectedly relevant to what users
mean when they say a model has answered the wrong question.

The translation to artificial systems began with the question
of whether a symbol-manipulating system could ever truly mean
anything at all. Stevan Harnad's 1990 paper \emph{The Symbol
Grounding Problem} asked how the symbols inside a computer
could come to refer to things in the world, given that the
symbols themselves are just shapes. The question became
urgent when statistical language models began producing
fluent text without any obvious mechanism for connecting that
text to anything outside the corpus.

Emily Bender and Alexander Koller's 2020 paper \emph{Climbing
towards NLU} sharpened the argument. They proposed a thought
experiment in which an octopus, intelligent but confined to
listening to undersea cables that carry human conversation,
learns to imitate the conversation perfectly without ever
encountering the world the conversation is about. Could the
octopus, on the basis of imitation alone, be said to
understand? Bender and Koller's answer was no: a system
trained only on linguistic form cannot, by that training
alone, learn meaning, because meaning is a relation between
form and the world that the form alone does not contain.
Modern Large Language Models are, in this analogy, the
octopus.

Gary Marcus, in a series of papers and books over the same
period, argued that the failures of modern LLMs are not bugs
to be fixed but the expected consequence of an architecture
that has no concept of reference. The argument is contested.
What is not contested is that the failure modes are real,
that they are characteristic, and that they are not going
away merely as a consequence of further scaling.

\paragraph{Where we stand.} The model this chapter develops has five axes, arranged from
the mechanical to the human.

\textbf{Structure.} Whether the output is well-formed for the
purpose at hand. A JSON object that parses; a grammatical
sentence; a piece of code that runs. The most mechanical of
the five axes, and the one on which modern models perform
best.

\textbf{Meaning.} Whether the output is internally coherent
and says what it appears to say. A paragraph whose sentences
agree with one another; an argument that does not contradict
its own premises; a chain of reasoning that holds together.
Detectable by careful reading without any appeal to the
world.

\textbf{Context.} Whether the output is appropriate for the
situation, the purpose, and the audience. The right tone, the
right level of detail, the right register, the right answer
to the question the user actually asked.

\textbf{Grounding.} Whether the output connects to verifiable
reality or has drifted into plausible-sounding fabrication.
The factual claims in the output, evaluated against the
world.

\textbf{Environment.} Whether the prompt itself contained
enough of what the user knows about the situation for the
question to be answerable at all. The most human of the
axes; the one on which the model is most at the mercy of the
person typing.

The five axes form a ladder. At the bottom, the model is at
its strongest and failures are rare. At the top, the model is
in a precise sense at the mercy of the human, and failures
are inevitable unless the human compensates.

The pedagogical use of the model is twofold. For users, it is
a diagnostic ladder: when an answer arrives that smells wrong,
walk the axes and ask which has failed. For engineers and
operators, it is a way of mapping the vocabulary of the
previous chapters onto a shape that non-specialists can hold.
Cross-entropy training is a structure-and-meaning mechanism
that makes no native attempt at grounding. Retrieval is a
grounding mechanism that does not by itself solve context.
Tool use is an environment-importing mechanism whose contract
is structural. The ladder makes these connections explicit.

The five axes map single-output failures. Four further failure
classes exist alongside them that the ladder cannot reach:
normative, adversarial, drift, and complacency. The chapter
introduces all of these before demonstrating the ladder in
practice.

\paragraph{What goes wrong.} The model itself has failure modes worth disclosing in
advance.

The first is the temptation to use it mechanically, as if the
diagnosis of an output were a checklist that always returns a
single answer. Real failures often span multiple axes. A
model that is wrong about a fact and presents the fact in the
wrong tone has failed on grounding and context simultaneously.
Pulling the two apart can be artificial.

The second is the temptation to treat the axes as orthogonal.
They are not. Failures of structure tend to drag meaning with
them; failures of context tend to drag grounding. The axes
are a vocabulary, not an algebra.

The third is the deeper mistake of treating the model as a
solution rather than as a diagnostic vocabulary. The model
does not, by itself, fix anything. It tells you which earlier
chapter to consult. The fixing is the work the previous
chapters described.

\section{The five axes, in plain language}
\addcontentsline{toc}{section}{The five axes, in plain language}

\subsection*{Structure: is the output well-formed?}

A model's output is structurally correct if it satisfies the
formal constraints of the format being requested. A piece of
JSON that parses. A grammatically complete sentence. A
Python script that imports without error. A SQL query that
the database accepts.

Structure is the most mechanical of the five axes and is, for
that reason, where modern Large Language Models perform best.
A structural failure is the kind a parser can catch. The
model produces a JSON-like blob with a missing comma, or
writes Python whose indentation contradicts the syntax rules,
or opens a markdown code fence and forgets to close it.
These failures are real but increasingly rare in the current
generation of frontier models, and most production systems
catch them automatically through what engineers call
\emph{constrained decoding}: forcing the model to produce
output that satisfies a known schema.

In diagnosis, structural failures are the easiest to identify.
They are also the easiest to be falsely reassured by: a
syntactically valid JSON object can still mean the wrong
thing, say the wrong thing, and refer to a thing that does
not exist. A practitioner who confirms structure and stops has
confirmed the least interesting property of the output.

\subsection*{Meaning: is the output internally coherent?}

A structurally correct output can still fail at the level of
meaning. The model produces a paragraph whose first sentence
contradicts its third. The model defines a term in one
section and then uses the term differently in the next. The
model's argument relies on a premise it has elsewhere denied.
These failures are detectable by careful reading, without any
appeal to the outside world. They are failures of internal
coherence.

The mechanism that produces meaning failures is the same
mechanism that produces fluent text. A language model trained
on next-token prediction (Chapter~1) optimises for the
\emph{plausibility} of the next token given the prefix.
Plausibility is a local property. A token can be plausible at
every position in a long sequence and still produce, at the
end of the sequence, a global contradiction. A single
forward pass through the model has no architectural verifier;
the model emits each token conditioned on the prefix without
checking the output as a whole against itself.

Recent training tries to narrow the gap. The reasoning models
that began appearing in 2024 (the OpenAI o1 and o3 lines,
DeepSeek's R1, Anthropic's extended-thinking modes) learn to
walk back contradictions in their own reasoning during
generation, in effect performing self-checking inside the
inference call. Outside the model, techniques sample multiple
completions and compare them, ask the model to critique and
revise its own output, or hand the output to a separate
verifier. None of these is a guarantee in the way a parser is
for structure. Each is a mitigation that trades cost for
reliability.

This is the axis that benchmarks tend to overstate. A model
that scores 85 out of 100 on a multiple-choice reading
comprehension test may nevertheless produce, in free
generation, paragraphs that contradict themselves in subtle
ways. Domain expertise tends to unmask this: the technical
reader who knows the territory catches inconsistencies that
the casual reader does not.

In diagnosis, the question to ask is not whether the output
sounds right but whether it \emph{is} right considered as a
self-contained artefact. Read it once for fluency, then again
for contradictions. When the cost of a meaning failure
justifies the additional compute, ask the model to check its
own work or sample multiple completions and compare them.

\subsection*{Context: is the output appropriate for the
situation?}

The third axis is whether the output is appropriate for its
situation, purpose, and audience. A model can produce a
structurally valid, internally coherent paragraph that
nevertheless answers the wrong question, addresses the wrong
audience, strikes the wrong tone, or pitches the wrong level
of detail. The technical literature calls this
\emph{pragmatic competence}; the user recognises it as the
model having ``answered a different question''.

Examples are easy to find. A user asks for a summary of a
research paper for a non-specialist audience and receives a
paragraph dense with jargon. A user asks a quick yes-or-no
question and receives a five-paragraph essay with the
yes-or-no buried in the middle. A user asks a clarifying
question about a sensitive topic and receives a generic
safety disclaimer instead of the requested clarification.
None of these are factually wrong. All of them are wrong for
the situation.

The mechanism is largely the gap between what the user knows
about the situation and what the prompt encodes. The model
has access only to the words in its context window. Anything
the user knows about themselves, their audience, their goals,
or their constraints that is not written into the prompt does
not exist for the model. Few-shot prompting (Chapter~6) is,
among other things, an attempt to encode situational
expectations through examples. System prompts (Chapter~12)
are an attempt to fix the situation up front. Alignment
training (Chapter~9) shapes the model's defaults along this
axis --- politeness, refusal, length, formality --- but the
defaults are defaults, not the situation.

In diagnosis, context failures are easy to recognise and
harder to fix. The user can usually point at the answer and
say what is wrong with it (``too long'', ``too technical'',
``misses the actual question''). The fix is almost always to
give the model more of the situation: who you are, who you
are writing for, why you need the answer, what would count as
a good one. The asymmetry between the ease of recognition and
the difficulty of fix is the diagnostic signature of a context
failure.

\subsection*{Grounding: does the output connect to reality?}

The fourth axis is whether the output connects to verifiable
reality or has drifted into plausible-sounding fabrication.
The practitioner literature calls a grounding failure a
\emph{hallucination}: the model produces a fluent, confident,
well-formed claim that does not correspond to any fact in
the world.

Grounding failures are uniquely dangerous because no internal
reading of the output reveals them. A hallucinated citation
has the same structural form as a real one. A hallucinated
statistic has the same surface plausibility as a true one. A
hallucinated historical fact reads with the same confidence
as a verified one. The check must come from outside the text.
The model has produced the most probable continuation of its
prompt; the most probable continuation is sometimes true,
sometimes false, and the model cannot, on its own, tell the
two apart.

The mechanism is the next-token prediction objective itself.
A model trained to predict the next token given the prefix
optimises for plausibility, not for truth. There is no signal
in the training that distinguishes a fact the model has
correctly memorised from a confabulation that fits the
surrounding pattern. This is the deeper philosophical point
that Harnad and Bender and Koller made: a system trained only
on form has no native way to ground form in reference.

The remedy at the system level is retrieval (Chapter~10). A
retrieval-augmented generation pipeline replaces the question
``what is the most probable thing to say'' with ``what is the
most probable thing to say given these specific sources''.
The grounding is no longer the model's problem; it is the
retriever's problem and the prompt's problem. Tool use
(Chapter~10 and Chapter~12) extends the same idea: a model
that can call a calculator can delegate arithmetic; a model
that can query a database can delegate fact lookup. The
pattern in each case is to take the work that demands
grounding and move it outside the language model entirely.

Grounding failures are compounded by a problem the model
cannot solve from the inside. The model writes with the same
assured, confident prose regardless of whether the underlying
claim is well-attested or thinly supported. A sentence about
a famous historical fact and a sentence about an invented one
look identical on the page. This means the reader cannot use
the model's apparent certainty as a guide to which claims
need checking. A hedged sentence does not mean the claim is
uncertain; a confident sentence does not mean it is reliable.
The practical consequence is that verification cannot be
applied selectively to claims that sound uncertain. All
factual claims need the same scrutiny --- not because all are
equally likely to be wrong, but because the surface of the
text gives no reliable signal of which are.

In diagnosis, grounding failures call for a particular
discipline. Treat any factual claim by the model as a
hypothesis. If the claim matters, verify it. If it cannot be
verified, do not use it. The careful practitioner builds
workflows that make verification cheap (retrieval, citations,
linked sources) so that grounding failures are catchable by
construction.

When grounding failures appear consistently in a specific
domain or type of question, the root cause may be a mismatch
between where the model was calibrated and where it is being
used. The per-call remedy --- better retrieval, more
verification --- addresses individual errors but not the
underlying pattern. The systemic remedy is evaluation on the
actual deployment population before relying on the tool at
scale.

\subsection*{Environment: did the prompt encode the situation?}

The fifth axis sits at the human end of the ladder, and is
the hardest to see because, by construction, it is invisible.
An environment failure occurs when the answer is wrong because
the prompt did not contain enough of the situation for the
question to be answerable. The user knows things --- about
themselves, their domain, their constraints, the local
conventions of their organisation, the events of yesterday
afternoon --- that the model does not. If the user does not
write those things into the prompt, the model cannot use
them. The model can produce a fluent, internally coherent,
grounded, well-formed answer that is wrong because the
question itself was underspecified.

This is the failure mode that experienced LLM users learn to
recognise first and explain to others last. It looks, from
the inside, like a model failure: the answer is wrong, so the
model must be wrong. From the outside, it is a prompt
failure: the model gave the most reasonable answer to the
question it was actually asked, which differs from the
question the user thought they were asking. The gap between
the two is the user's environment.

The mechanism is fundamental and not fixable by training. A
language model has access to its weights and its context
window. It does not have access to the user's calendar, the
user's team, the user's organisation's coding conventions,
the user's medical history, the user's project's constraints,
or any of the several hundred other things that shape what
counts as a good answer in a real situation. Some of these
can be imported into the context: through retrieval, through
tool use, through careful prompting, through long-running
agent contexts that accumulate state. None of them can be
imported automatically.

In diagnosis, an environment failure has a characteristic
shape. The user reads the answer and feels frustrated; the
answer is, on its own terms, fine; what is wrong is that it
does not address the real situation. The fix is rarely to ask
the model again. The fix is to write down, explicitly, the
parts of the environment the model needed to see.
Practitioners who become good at LLMs are practitioners who
have learned which parts of their environment need importing
for which kinds of question.

\section{What the ladder does not catch}
\addcontentsline{toc}{section}{What the ladder does not catch}

The five axes map the full range of single-output failures:
everything that can go wrong when a single prompt produces a
single response. But several failure classes exist alongside
them that fall outside that single-output frame entirely.
These are worth knowing before walking the ladder, because
they require different detection strategies and different
remedies --- and the ladder, by itself, does not surface them.

The first is the \emph{normative} failure: an output that is
structurally fine, internally coherent, situationally
appropriate, factually grounded, and environmentally informed
but that is, by some standard the user holds, simply wrong
--- offensive, biased, unethical, in conflict with the values
the user expected the system to hold. These failures live in
Chapter~9's territory rather than this one. The five axes
ask whether the answer is correct; the normative dimension
asks whether the answer is right, and right is not the same
thing as correct. A related version operates at the
population level: an output correct for one user may be part
of a pattern that, across many users, produces systematic
harm. Chapter~11 is where population-level monitoring belongs.

The second is the \emph{adversarial} failure: an output
induced by a malicious party manipulating one of the inputs
--- the prompt, a retrieved document, a tool result. The
ladder will identify the failure (most often as a context or
grounding failure), but the diagnosis does not explain that
an attacker caused it. A particularly subtle form: an
attacker who can place content in documents the model will be
asked to read --- a contract being reviewed, a web page being
summarised, an email being processed --- can embed
instructions in that content. The model, trained to follow
instructions, may follow them without signalling that
anything is amiss. Chapter~12 is where the discipline of
prompt injection and confused-deputy defences belongs.

The third is the \emph{drift} failure: an output that is
fine considered as a single artefact but that, considered as
part of a stream of similar outputs over time, indicates that
the system's behaviour has changed in a way the operator did
not anticipate. Drift is detected statistically rather than
individually, and the monitoring infrastructure of
Chapter~11 is what catches it.

The fourth is the \emph{complacency} failure. The ladder
assumes an active reviewer examining each output carefully.
Human-factors research has documented what happens to
vigilance in high-volume, mostly-correct review settings:
after a long run of correct outputs, attention decays. The
reviewer is still present; the review becomes perfunctory.
The record shows that review happened; it does not show how
carefully. The effective error rate of a human-plus-model
system can rise substantially over time not because the model
degrades but because the reviewer does. The remedy is in the
workflow: rotation, sampling-based spot checks, and review
metrics that measure quality of review rather than merely its
occurrence.

The four classes require different detection strategies than
the ladder provides. Normative failures call for governance
and ethics review. Adversarial failures call for pipeline
defences and adversarial testing. Drift calls for statistical
monitoring. Complacency calls for workflow design. The ladder,
which follows, addresses the five-axis single-output failures
directly. Knowing both the map and the tool makes the walk
more useful.

\section{Walking the ladder}
\addcontentsline{toc}{section}{Walking the ladder}

With the full failure landscape in view, the ladder's scope
is clear: it is a diagnostic procedure for the five-axis
failures --- the single-output failures that a systematic
walk can surface. Walk it bottom-up.

\begin{enumerate}
  \item \textbf{Structure.} Does the output parse, compile,
        validate against its expected form? If not, the
        failure is structural. The remedy lives in the model's
        decoding (Chapter~6) or the system's constrained
        generation (Chapter~12).
  \item \textbf{Meaning.} Read the output as a self-contained
        artefact. Does it contradict itself? Does it use a
        term differently in different sentences? Does its
        argument rely on a premise it has elsewhere denied?
        If yes, the failure is one of meaning. The remedies
        are sampling multiple completions and comparing them,
        asking the model to check its own work, or using
        chain-of-thought to externalise the reasoning.
  \item \textbf{Context.} Is the output appropriate for the
        situation --- the right tone, the right level, the
        right format, the right answer to the actual question?
        If not, the failure is contextual. The remedy is to
        import more of the situation into the prompt: who you
        are, who you are writing for, why you need the answer,
        what would count as a good one.
  \item \textbf{Grounding.} Are the factual claims in the
        output true? If verification is hard, treat the claims
        as hypotheses. If verification is easy, do it.
        Failures here call for retrieval, tool use, or human
        review.
  \item \textbf{Environment.} Is the question itself fully
        specified? Did the prompt contain enough of the user's
        situation for the question to have a single right
        answer? If not, no amount of model improvement will
        help; the user must supply more of their environment.
\end{enumerate}

The discipline of walking the ladder bottom-up is that
lower-axis failures are typically easier to diagnose and
cheaper to remedy than higher-axis ones. A practitioner who
jumps straight to ``the model hallucinated'' may have missed
a meaning or context failure that does not require any change
to the model at all. Conversely, a practitioner who never
considers grounding will deploy systems that are confidently
wrong about the world.

When a single output fails on more than one axis, the
procedure is to label all axes that fail and to address them
in ladder order: the lowest-axis failure is fixed first,
because higher-axis fixes assume the lower axes hold. A useful
heuristic for prioritising the work is to ask which axis, if
it had been correct, would have made the others correctable.
A grounded answer in the wrong tone can be fixed by the user
through a follow-up; a fluent answer about a fact that does
not exist cannot be fixed by the user at all, because the
user did not know the fact was wrong. The limiting axis is
the one to pay first.

For workflows that chain multiple steps --- summarising
several documents, then comparing the summaries, then drawing
a conclusion --- apply the ladder at each intermediate output,
not only at the final result. A chain where every individual
step passes the ladder is not therefore safe: errors from an
early step enter the next as premises and compound.

\section{A worked diagnosis}
\addcontentsline{toc}{section}{A worked diagnosis}

The ladder is easier to use after you have walked it once.
Here is a short worked example: a plausible, confident, fluent
response that contains a real failure, diagnosed axis by axis.

\medskip
\noindent\textbf{The prompt:}
\begin{quote}
\emph{Can you tell me who wrote the novel \textit{1984} and give
me a brief summary of the plot?}
\end{quote}

\noindent\textbf{The model's response:}
\begin{quote}
\emph{1984 was written by George Orwell and published in 1949.
The novel is set in a totalitarian society called Oceania,
ruled by an organisation known as the Party and its figurehead
Big Brother. The protagonist Winston Smith works at the Ministry
of Truth, where his job is to rewrite historical records to
match the Party's current version of events. Winston begins a
forbidden love affair with a woman named Julia, and together
they contact a man called O'Brien, who they believe leads a
secret resistance movement. The novel culminates in Winston
being captured and brought to Room 101 in the Ministry of
Truth, where he is psychologically broken by being confronted
with his greatest fear. The novel's title, many believe, was
formed by reversing the last two digits of 1948, the year
Orwell was writing it.}
\end{quote}

\medskip
\noindent Walk the ladder.

\medskip
\textbf{Structure.} Grammatically correct sentences, clear
paragraphs, no malformed output. Structure passes.

\textbf{Meaning.} The response is internally coherent. Each
sentence follows naturally from the last. There are no
contradictions. The argument about Winston's arc holds
together as written. Meaning passes.

\textbf{Context.} The tone is appropriate for a general
knowledge question. The length is reasonable. The response
answers the actual question asked. Context passes.

\textbf{Grounding.} Here the ladder catches something.
Room 101 is not in the Ministry of Truth. In the novel it
is in the Ministry of Love --- the building responsible for
torture and psychological reconditioning. Winston works in
the Ministry of Truth, but he is taken to a different
building entirely for his interrogation. The two ministries
are distinct. The model's response conflates them in a
single sentence that sounds correct because it has the
right surface shape (a place, an action, a consequence)
even though the detail is wrong.

The final sentence is a different kind of grounding problem.
The claim that Orwell arrived at the title 1984 by reversing
the digits of 1948 is plausible and widely repeated.
Orwell himself never confirmed it in any letter or diary
entry that has come to light. It may be true; it may be a
story that accumulated credibility by being repeated. The
model presents it as common knowledge (``many believe'')
rather than as a verified fact --- a small hedge, but the
hedge is easy to miss. A reader who passed the claim along
would be passing along an unverified attribution with a
degree of confidence it does not earn.

\textbf{Environment.} The question is specific and clear.
No underspecification. Environment passes.

\medskip
\noindent\textbf{Diagnosis:} one grounding failure
(wrong ministry for Room 101) and one disguised grounding
problem (an unverified attribution presented as common
knowledge). The fix is not to distrust the model's account of
the novel's themes or characters, which is accurate. The fix
is to treat any specific, checkable detail --- a place, a date,
an attribution --- as a hypothesis, and to verify it if the
answer will be used somewhere it matters.

\medskip
This is the discipline the ladder recommends. Structure,
meaning, and context passed; the response reads fluently; the
novel's themes are described correctly; the failure is
concentrated in two specific details that only a reader who
already knows the novel, or who checks, will catch. That
concentration is characteristic of grounding failures. They
are not randomly distributed. They tend to appear in the
peripheral details --- the supporting citations, the secondary
attributions, the names and dates that surround a well-known
core --- precisely because those details are statistically
underrepresented in training and the model interpolates them
from pattern rather than from memory.

\section{Connections to earlier chapters}
\addcontentsline{toc}{section}{Connections to earlier chapters}

The five axes connect to the technical mechanisms of earlier
chapters in ways worth making explicit.

\textbf{Structure} is the territory of the decoding step
(Chapter~6) and the structured-output contracts of tool use
(Chapter~12). When you see a structural failure, the
intervention is in the decoder, the schema, or the validation
layer immediately after the model.

\textbf{Meaning} is the territory of the next-token prediction
objective (Chapter~1), the deep neural network architecture
(Chapter~3), and the inference-time techniques that trade
compute for coherence (Chapter~6). When you see a meaning
failure, the intervention is in how the model is being asked
to think out loud, how many samples are being drawn, and
whether self-checking is being used.

\textbf{Context} is the territory of in-context learning
(Chapter~6), the attention mechanism (Chapter~5), the
prompt-construction discipline (Chapter~10), the system
prompts of the API surface (Chapter~12), and the alignment
work that shapes the model's defaults (Chapter~9). When you
see a context failure, the intervention is in the prompt: more
of the situation, more explicit constraints, more examples of
the desired behaviour.

\textbf{Grounding} is the territory of retrieval (Chapter~10),
tool use, the symbol-grounding literature (Harnad, Bender and
Koller), and the memorisation-versus-hallucination distinction
of Chapter~6. When you see a grounding failure, the
intervention is to move the work that demands ground truth
outside the model: into a retriever, a tool, a verification
step, a human reviewer.

\textbf{Environment} is the territory of system design rather
than model design. Long-context architectures (Chapter~5)
extend how much environment can be imported. Agent loops
(Chapter~12) accumulate environment over time. The
documentation, logging, and incident reporting disciplines of
governance (Chapter~11) make the environment the system sees
explicit and auditable.

The general pattern is that lower axes are increasingly the
\emph{model's} responsibility, while higher axes are
increasingly the \emph{system's} responsibility, and at the
top the responsibility belongs to the human writing the
prompt. The ladder makes the locus of intervention explicit.


\section*{Closing thought}
\addcontentsline{toc}{section}{Closing thought}

You have, by this point, walked through the entire engineering
of a Large Language Model in plain language. You know,
roughly, where they came from, what they do, why they do it
the way they do, and how they fail. The vocabulary you now
have --- structure, meaning, context, grounding, environment;
hallucination, sycophancy, jailbreak, prompt injection,
alignment tax; pretraining, alignment, retrieval, tool
use; Markov, Shannon, Rumelhart, Mikolov, Vaswani, Hu,
Hoffmann, Christiano, Bender --- is the vocabulary of a
literate observer of the field as it stands at the time of
this writing. It will age. The technology will change. The
deeper structure of why models do what they do, and why they
fail in the ways they do, will not.

The most useful posture, in our experience, is the one we
recommended in the preface and have tried to model throughout.
These systems are remarkable, in specific and definable ways.
They are limited, in specific and definable ways. The public
conversation has tended to oscillate between the two halves
of that observation. The reader who has come this far knows
better. The model is doing exactly what we asked. What we
asked is not quite what we meant. The space between the two is
the field. The vocabulary you now have is what makes that
space inhabitable.

\section*{Bridges}
\addcontentsline{toc}{section}{Bridges}

This chapter built on every preceding chapter. It is the
chapter that the others were written to make possible. The
five-axis ladder is not a fix for any of the failure modes
catalogued in the previous twelve chapters; it is a way of
holding all of them at once, in language a non-specialist can
use, so that a wrong answer in front of you can be diagnosed
and acted on rather than merely lamented.

There are no further chapters. The remainder of the work is
yours: to use these systems thoughtfully, to ask the
questions of vendors and engineers and regulators that the
preceding chapters have prepared you to ask, and to keep an
honest record of what you find. The next edition of this
guide, if there is one, will be better for what you find than
for what its authors imagined.
